\chapter{Manjul Bhargava --- Fields Medal 2014}\label{ch2014manjulbhargava:manjul-bhargava-fields-medal-2014-a-comprehensive-report-on-his-work-and-impact}

\begin{center}\rule{0.5\linewidth}{0.5pt}\end{center}

\textbf{Prepared:} July 2026 \textbf{Subject:} Manjul Bhargava (b.
August 8, 1974, Hamilton, Ontario, Canada) \textbf{Position:} Robert C.
Gunning and R. Brandon Fradd Professor of Mathematics, Princeton
University \textbf{Doctoral Advisor:} Andrew Wiles \textbf{Fields Medal
Citation:} \emph{``For developing powerful new methods in the geometry
of numbers, which he applied to count rings of small rank and to bound
the average rank of elliptic curves.''}

\begin{center}\rule{0.5\linewidth}{0.5pt}\end{center}

\section{Table of Contents}\label{ch2014manjulbhargava:table-of-contents}

\begin{enumerate}
\def\labelenumi{\arabic{enumi}.}
\tightlist
\item
 Introduction
\item
 Biographical Background and Education
\item
 The Geometry of Numbers: Historical Context
\item
 Gauss Composition and the Higher Composition Laws
\item
 The Bhargava Cube: A New View of Gauss Composition
\item
 Higher Composition Laws II: Cubic Analogues
\item
 Higher Composition Laws III: Parametrization of Quartic Rings
\item
 Higher Composition Laws IV: Quintic Rings
\item
 The Density of Discriminants of Number Fields
\item
 The Bhargava Factorial
\item
 The Average Rank of Elliptic Curves
\item
 The Bhargava--Shankar Program on Selmer Groups
\item
 The 2-Selmer Group and Binary Quartic Forms
\item
 Ternary Cubic Forms and the 3-Selmer Group
\item
 The 4-Selmer and 5-Selmer Groups
\item
 Consequences for the BSD Conjecture and Rank Distribution
\item
 Universal Quadratic Forms and the 290-Theorem
\item
 Further Contributions: Hyperelliptic Curves, Plane Cubics, and the
 Hasse Principle
\item
 Broader Impact and Influence
\item
 Awards and Recognition
\item
 Selected Major Publications
\item
 Ongoing Work and Future Directions
\item
 Conclusion
\item
 References and Further Reading
\end{enumerate}

\begin{prereqs} Elementary number theory (quadratic forms, class numbers), abstract algebra (rings of integers, group actions), basic algebraic geometry. \end{prereqs}

\begin{objectives} Explain higher composition laws and the Bhargava boxes; describe the average size of 2-Selmer groups; understand the connection between number fields and geometry. \end{objectives}


\begin{center}\rule{0.5\linewidth}{0.5pt}\end{center}

\section{1. Introduction}\label{ch2014manjulbhargava:introduction}

Manjul Bhargava was awarded the Fields Medal in 2014 at the age of 39
for a series of spectacular breakthroughs in number theory. His work
spans the geometry of numbers, arithmetic statistics, algebraic number
theory, and the arithmetic of elliptic curves. In each of these areas,
he has solved problems that had resisted progress for decades, often by
discovering unexpected connections between classical mathematics and
modern techniques.

Bhargava's career is notable for its precocity. As a graduate student
at Princeton under Andrew Wiles, he discovered 14 new composition laws
generalizing Gauss's classical law for binary quadratic forms --- a
discovery that came to him while playing with a Rubik's cube. This work,
completed as his doctoral thesis, was immediately recognized as a major
advance and published in the Annals of Mathematics in a series of four
papers. In the years that followed, he developed powerful new methods
for counting number fields of small degree, determined the density of
discriminants of quartic and quintic fields (problems that had been open
for over a century), and, in a landmark series of collaborations with
Arul Shankar, proved that the average rank of elliptic curves is bounded
--- a result that revolutionized the arithmetic statistics of elliptic
curves.

Bhargava's contributions are characterized by their combination of deep
classical learning with striking originality. He is as comfortable with
Gauss's Disquisitiones Arithmeticae as with the most modern algebraic
geometry, and he has a gift for finding simple, elegant formulations of
problems that had seemed intractable. His work on higher composition
laws showed that Gauss's composition --- long thought to be a unique
phenomenon for quadratic forms --- was just the first example of a
general family of laws. His work on counting number fields brought the
power of the geometry of numbers to problems that had been studied since
the time of Dirichlet. And his work on elliptic curves transformed a
subject that had been dominated by deep conjectures into one with
precise, provable theorems.

This report provides a detailed account of Bhargava's important
contributions, the mathematical context in which they arose, the
techniques he developed, and their lasting impact on number theory.

\begin{center}\rule{0.5\linewidth}{0.5pt}\end{center}

\section{2. Biographical Background and
Education}\label{ch2014manjulbhargava:biographical-background-and-education}

\subsection{2.1 Early Life}\label{ch2014manjulbhargava:early-life}

Manjul Bhargava was born on August 8, 1974, in Hamilton, Ontario,
Canada, to parents of Indian origin. His mother, Mira Bhargava, is a
mathematician who taught at Hofstra University and was his first
mathematics teacher. His father came from the Jaipur region of
Rajasthan, India. The family moved to the United States when Bhargava
was young, and he grew up primarily on Long Island, New York.

From an early age, Bhargava showed exceptional mathematical talent. His
mother taught him arithmetic and basic mathematics at home, and he
rapidly outpaced his peers. He also developed a deep interest in
Sanskrit and Indian classical music, particularly the tabla, a
traditional Indian drum. He would later study tabla under the notable
Zakir Hussain.

\subsection{2.2 Education}\label{ch2014manjulbhargava:education}

Bhargava attended high school on Long Island, completing advanced
mathematics courses by age 14. He then entered Harvard University, where
he earned his A.B. in mathematics summa cum laude in 1996. As an
undergraduate, he wrote a thesis on what would become the Bhargava
factorial, a generalization of the classical factorial function. This
work won him the Morgan Prize for outstanding undergraduate research in
mathematics.

After Harvard, Bhargava moved to Princeton University for his doctoral
studies. His PhD advisor was Andrew Wiles, who had famously proved
Fermat's Last Theorem just a few years earlier. Under Wiles's
supervision, Bhargava wrote a thesis on higher composition laws,
discovering 14 new laws generalizing Gauss's composition. He received
his PhD in 2001.

\subsection{2.3 Academic Career}\label{ch2014manjulbhargava:academic-career}

After a year as a visiting scholar at the Institute for Advanced Study
(2001--2002) and a year at Harvard (2002--2003), Bhargava joined the
Princeton faculty as a tenured full professor in 2003 at the age of 28
--- one of the youngest tenured professors in the university's history.
He has remained at Princeton ever since, currently holding the Robert C.
Gunning and R. Brandon Fradd Professorship of Mathematics. He also holds
the Stieltjes Chair in Number Theory at Leiden University in the
Netherlands.

Bhargava is known for his dedication to teaching and public engagement.
He teaches a popular freshman seminar called ``The Mathematics of Magic
Tricks and Games'' and has served as the inaugural Distinguished Chair
for Public Dissemination of Mathematics at the National Museum of
Mathematics in New York.

\begin{center}\rule{0.5\linewidth}{0.5pt}\end{center}

\section{3. The Geometry of Numbers: Historical
Context}\label{ch2014manjulbhargava:the-geometry-of-numbers-historical-context}

\subsection{3.1 Minkowski's Legacy}\label{ch2014manjulbhargava:minkowskis-legacy}

The geometry of numbers was founded by Hermann Minkowski in the late
19th century. Its central idea is to use geometric methods ---
particularly the study of lattices and convex bodies --- to solve
problems in number theory. Minkowski's fundamental theorem states that a
centrally symmetric convex body of volume greater than \(2^n\) times the
covolume of a lattice must contain a non-zero lattice point.

This simple idea has significant consequences. It gives a geometric proof
of the Dirichlet unit theorem, provides bounds on the discriminant of
number fields, and forms the basis for much of modern algebraic number
theory.

\subsection{3.2 The Davenport--Heilbronn
Program}\label{ch2014manjulbhargava:the-davenportheilbronn-program}

In the 1960s and 1970s, Harold Davenport and Hans Heilbronn developed a
powerful extension of Minkowski's ideas. They showed how to count the
number of cubic fields with bounded discriminant by counting integer
orbits of a certain group action. Their work gave the first asymptotic
formula for cubic fields:

\[
N_3^{\pm}(X) \sim \frac{1}{3\zeta(3)} X \quad \text{and} \quad \frac{1}{4\zeta(3)} X
\]

for totally real and complex cubic fields respectively.

This result was a landmark in arithmetic statistics, but it was limited
to the cubic case. For degrees 4 and 5, the problem remained open for
decades. The obstruction was that the necessary parametrizations of
quartic and quintic rings were unknown.

\subsection{3.3 Bhargava's Entry}\label{ch2014manjulbhargava:bhargavas-entry}

Bhargava's higher composition laws provided exactly the parametrizations
needed to extend the Davenport--Heilbronn program to higher degrees. His
work showed that quartic rings could be parametrized by pairs of ternary
quadratic forms, and quintic rings by quadruples of alternating 5-by-5
bilinear forms. With these parametrizations in hand, Bhargava was able
to count quartic and quintic fields asymptotically, solving problems
that had been open for over 50 years.

\begin{center}\rule{0.5\linewidth}{0.5pt}\end{center}

\section{4. Gauss Composition and the Higher Composition
Laws}\label{ch2014manjulbhargava:gauss-composition-and-the-higher-composition-laws}

\subsection{4.1 Gauss's Original Composition
Law}\label{ch2014manjulbhargava:gausss-original-composition-law}

In his 1801 masterwork \emph{Disquisitiones Arithmeticae}, Gauss
discovered a notable law of composition for binary quadratic forms. A
binary quadratic form is an expression of the form:

\[
Q(x,y) = ax^2 + bxy + cy^2
\]

where \(a, b, c\) are integers. Gauss showed that one can ``compose''
two such forms to obtain a third form, and that this operation endows
the set of equivalence classes of forms of a given discriminant with a
group structure.

This group is isomorphic to the ideal class group of the quadratic
number field \(\mathbb{Q}(\sqrt{D})\), a connection that foreshadowed
the development of algebraic number theory by Dedekind and others.

Gauss's proof of the composition law was notoriously complicated: it
took 20 pages of dense calculations. For 200 years, this was considered
the definitive treatment.

\subsection{4.2 The Search for Higher Composition
Laws}\label{ch2014manjulbhargava:the-search-for-higher-composition-laws}

After Gauss, mathematicians wondered whether similar composition laws
existed for forms of higher degree. Binary cubic forms, ternary
quadratic forms, and higher-degree forms were all studied, but no
general composition laws were found.

The prevailing view became that Gauss composition was a unique
phenomenon, a consequence of the special structure of binary quadratic
forms. This view persisted for two centuries.

\subsection{4.3 Bhargava's Breakthrough}\label{ch2014manjulbhargava:bhargavas-breakthrough}

As a graduate student, Bhargava read the \emph{Disquisitiones} and was
struck by the complexity of Gauss's proof. He became convinced that
there had to be a simpler way. The breakthrough came in the summer of
1998 while he was playing with a 2×2×2 Rubik's cube.

The insight was to label the eight corners of a cube with integers and
then consider the three ways of slicing the cube into two 2×2 layers.
Each slicing produces a pair of 2×2 matrices, and from each pair one can
construct a binary quadratic form. The three forms obtained from the
three slicings have the same discriminant, and the composition of any
two yields the third --- exactly reproducing Gauss's composition law.

This geometric picture not only simplified Gauss's proof but revealed
that Gauss composition was just one instance of a much more general
phenomenon. By considering cubes of varying sizes and dimensions,
Bhargava discovered 14 new composition laws, including laws for cubic
forms, quartic forms, and forms in more variables.

\begin{center}\rule{0.5\linewidth}{0.5pt}\end{center}

\section{5. The Bhargava Cube: A New View of Gauss
Composition}\label{ch2014manjulbhargava:the-bhargava-cube-a-new-view-of-gauss-composition}

\subsection{5.1 The Cube Construction}\label{ch2014manjulbhargava:the-cube-construction}

A 2×2×2 cube of integers \(a_{ijk}\) has eight entries. Slice the cube
along each of its three axes to obtain three pairs of 2×2 matrices:

\begin{itemize}
\tightlist
\item
 Front/back: \((M_1, N_1)\)
\item
 Left/right: \((M_2, N_2)\)
\item
 Top/bottom: \((M_3, N_3)\)
\end{itemize}

From each pair \((M, N)\), one constructs a binary quadratic form:

\[
Q_i(x, y) = -\det(M_i x + N_i y)
\]

\subsection{5.2 The Composition Law}\label{ch2014manjulbhargava:the-composition-law}

Bhargava proved that the three forms \(Q_1, Q_2, Q_3\) all have the same
discriminant, and that the composition of any two (in the sense of
Gauss) is equivalent to the third:

\[
Q_1 \circ Q_2 = Q_3
\]

up to equivalence. This gives a remarkably simple proof of Gauss's
composition law.

\subsection{5.3 The Group Law on Elliptic
Curves}\label{ch2014manjulbhargava:the-group-law-on-elliptic-curves}

Bhargava also observed that the cube composition resembles the group law
on an elliptic curve. Given three points \(P, Q, R\) on an elliptic
curve with \(P + Q + R = 0\), the coordinates of the points satisfy
relations analogous to the cube composition. This observation connected
Gauss composition to one of the central objects of modern number theory.

\subsection{5.4 Generalizations}\label{ch2014manjulbhargava:generalizations}

The cube construction admits natural generalizations:

\begin{itemize}
\tightlist
\item
 \textbf{3×3×3 cubes} yield composition laws for cubic forms.
\item
 \textbf{2×2×2×2 hypercubes} yield composition laws for forms of higher
 degree.
\item
 \textbf{Cubes of integers modulo group actions} yield parametrizations
 of algebraic structures.
\end{itemize}

In total, Bhargava discovered 14 higher composition laws, published in a
well-known series of Annals of Mathematics papers.

\begin{center}\rule{0.5\linewidth}{0.5pt}\end{center}

\section{6. Higher Composition Laws II: Cubic Analogues of Gauss
Composition}\label{ch2014manjulbhargava:higher-composition-laws-ii-cubic-analogues-of-gauss-composition}

\subsection{6.1 The Cubic Case}\label{ch2014manjulbhargava:the-cubic-case}

After establishing the cube interpretation of Gauss composition,
Bhargava turned to higher-degree analogues. The first and most natural
generalization is to cubic forms.

A binary cubic form is an expression:

\[
C(x, y) = ax^3 + bx^2 y + cxy^2 + dy^3
\]

The question: is there a composition law for cubic forms analogous to
Gauss's law for quadratic forms?

\subsection{6.2 The 3×3×3 Cube}\label{ch2014manjulbhargava:the-333-cube}

Bhargava showed that a 3×3×3 cube of integers gives rise to three binary
cubic forms, one from each slicing. These three forms have the same
discriminant, and the composition of any two yields the third. This is
the cubic analogue of Gauss composition.

\subsection{6.3 Connections to Class
Groups}\label{ch2014manjulbhargava:connections-to-class-groups}

Just as Gauss composition for quadratic forms relates to class groups of
quadratic fields, the cubic composition law relates to class groups of
cubic fields. Bhargava used this connection to study the distribution of
cubic fields and to compute averages of class group invariants.

\begin{center}\rule{0.5\linewidth}{0.5pt}\end{center}

\section{7. Higher Composition Laws III: Parametrization of Quartic
Rings}\label{ch2014manjulbhargava:higher-composition-laws-iii-parametrization-of-quartic-rings}

\subsection{7.1 The Problem}\label{ch2014manjulbhargava:the-problem}

A quartic ring is a commutative ring with identity that is free of rank
4 as a \(\mathbb{Z}\)-module. Understanding the set of all quartic rings
--- and especially those that are maximal orders in quartic number
fields --- is a fundamental problem in algebraic number theory.

The Delone--Faddeev correspondence (1940s) gave a beautiful
parametrization of cubic rings by binary cubic forms. This was the key
to the Davenport--Heilbronn counting of cubic fields. For quartic rings,
no such parametrization was known.

\subsection{7.2 Bhargava's
Parametrization}\label{ch2014manjulbhargava:bhargavas-parametrization}

Bhargava discovered that quartic rings are parametrized by pairs of
ternary quadratic forms \((A, B)\) modulo the action of
\(\mathrm{GL}_2(\mathbb{Z}) \times \mathrm{SL}_3(\mathbb{Z})\). More
precisely, there is a bijection between:

\begin{itemize}
\tightlist
\item
 \(\mathrm{GL}_2(\mathbb{Z}) \times \mathrm{SL}_3(\mathbb{Z})\)-orbits
 of pairs of ternary quadratic forms with a non-degeneracy condition
\item
 Isomorphism classes of quartic rings
\end{itemize}

This is a higher composition law: the pair of ternary quadratic forms
can be ``composed'' to produce a quartic ring.

\subsection{7.3 The Invariants}\label{ch2014manjulbhargava:the-invariants}

For a pair of ternary quadratic forms \((A, B)\), Bhargava defined
invariants that correspond to the discriminant and other invariants of
the associated quartic ring. The discriminant of the quartic ring is
given by a polynomial in the coefficients of \(A\) and \(B\).

\subsection{7.4 Impact}\label{ch2014manjulbhargava:impact}

This parametrization was the crucial ingredient for counting quartic
fields. Without it, the asymptotic enumeration of quartic fields was
impossible. With it, Bhargava was able to prove that the number of
\(S_4\)-quartic fields with absolute discriminant at most \(X\) is
asymptotic to a constant times \(X\).

\begin{center}\rule{0.5\linewidth}{0.5pt}\end{center}

\section{8. Higher Composition Laws IV: Quintic
Rings}\label{ch2014manjulbhargava:higher-composition-laws-iv-quintic-rings}

\subsection{8.1 Quintic Rings}\label{ch2014manjulbhargava:quintic-rings}

Paralleling the quartic case, Bhargava found a parametrization of
quintic rings by quadruples of alternating 5×5 bilinear forms modulo the
action of
\(\mathrm{GL}_4(\mathbb{Z}) \times \mathrm{SL}_5(\mathbb{Z})\).

\subsection{8.2 The Parametrization
Theorem}\label{ch2014manjulbhargava:the-parametrization-theorem}

Specifically, Bhargava proved a bijection between:

\begin{itemize}
\tightlist
\item
 \(\mathrm{GL}_4(\mathbb{Z}) \times \mathrm{SL}_5(\mathbb{Z})\)-orbits
 of quadruples of alternating 5×5 bilinear forms with a non-degeneracy
 condition
\item
 Isomorphism classes of quintic rings
\end{itemize}

\subsection{8.3 Counting Quintic Fields}\label{ch2014manjulbhargava:counting-quintic-fields}

With this parametrization, Bhargava proved the asymptotic formula:

\[
N_5^{\pm}(X) \sim C_5^{\pm} X
\]

where \(N_5^{\pm}(X)\) counts the number of quintic fields with
signature \((r_1, r_2)\) (totally real or totally complex) and absolute
discriminant at most \(X\), and \(C_5^{\pm}\) are explicit constants.

This was the first asymptotic result for number fields of degree 5, a
problem that had been open since the time of Dedekind. It showed that a
positive proportion of quintic fields have Galois group \(S_5\) (the
full symmetric group), and that the remaining quintic fields are
exponentially rarer.

\subsection{8.4 The Power of Higher Composition
Laws}\label{ch2014manjulbhargava:the-power-of-higher-composition-laws}

The series of papers on higher composition laws represents one of the
most impressive doctoral theses in the history of mathematics. In four
papers, each published in the Annals of Mathematics, Bhargava: -
Simplified and unified Gauss's composition law - Discovered 14 new
composition laws - Parametrized cubic, quartic, and quintic rings -
Opened the door to the arithmetic statistics of number fields

\begin{center}\rule{0.5\linewidth}{0.5pt}\end{center}

\section{9. The Density of Discriminants of Number
Fields}\label{ch2014manjulbhargava:the-density-of-discriminants-of-number-fields}

\subsection{9.1 The Classical Problem}\label{ch2014manjulbhargava:the-classical-problem}

The discriminant of a number field is one of its important
invariants. It measures the ramification of the field and appears in the
functional equation of the Dedekind zeta function. A central problem in
arithmetic statistics, going back to Gauss, is to understand the
distribution of discriminants of number fields.

Let \(N_n^{\pm}(X)\) denote the number of number fields of degree \(n\)
with signature \((r_1, r_2)\) and absolute discriminant at most \(X\).
The Minkowski bound implies that \(N_n^{\pm}(X)\) grows at most like
\(X^{n/2}\), but precise asymptotics were known only for \(n = 2\)
(quadratic fields, by Gauss) and \(n = 3\) (cubic fields, by
Davenport--Heilbronn).

\subsection{9.2 Bhargava's Asymptotics for Quartic
Fields}\label{ch2014manjulbhargava:bhargavas-asymptotics-for-quartic-fields}

Using his parametrization of quartic rings, Bhargava proved:

\[
N_4(X) \sim \frac{5}{24} \prod_p \left(1 + p^{-2} - p^{-3} - p^{-4}\right) \cdot X
\]

for \(S_4\)-quartic fields. The constant involves a product over all
primes, reflecting the local structure of quartic rings.

\subsection{9.3 Asymptotics for Quintic
Fields}\label{ch2014manjulbhargava:asymptotics-for-quintic-fields}

For quintic fields, Bhargava proved:

\[
N_5(X) \sim C_5 X
\]

where \(C_5\) is an explicit constant approximately equal to 0.0585 for
totally real quintics and 0.9363 for totally complex quintics. These
constants were the first --- and remain the only --- asymptotic
constants for number fields of degree 5.

\subsection{9.4 The Mass Formula for Local
Fields}\label{ch2014manjulbhargava:the-mass-formula-for-local-fields}

In addition to these counting results, Bhargava proved a mass formula
for extensions of local fields. This formula, which counts the number of
extensions of a local field with given ramification data, had been
conjectured by Serre and others. Bhargava's proof used the geometry of
numbers in a novel way, applying his parametrizations to the local
setting.

\subsection{9.5 The Cohen--Lenstra
Heuristics}\label{ch2014manjulbhargava:the-cohenlenstra-heuristics}

Bhargava's counting results have deep connections to the Cohen--Lenstra
heuristics, which predict the distribution of class groups of number
fields. His work on the average size of class groups of cubic fields
confirmed predictions of the Cohen--Lenstra heuristics in certain cases
and provided new evidence for their general validity.

\begin{center}\rule{0.5\linewidth}{0.5pt}\end{center}

\section{10. The Bhargava Factorial}\label{ch2014manjulbhargava:the-bhargava-factorial}

\subsection{10.1 The Definition}\label{ch2014manjulbhargava:the-definition}

As an undergraduate at Harvard, Bhargava introduced a notable
generalization of the factorial function. Given an infinite subset \(S\)
of the integers (or more generally, a Dedekind domain), Bhargava defined
a \(p\)-ordering of \(S\) for each prime \(p\) and used it to define
\(k!_S\), the Bhargava factorial of \(k\) with respect to \(S\).

The construction is as follows: - For a prime \(p\), choose a sequence
\(a_0, a_1, a_2, \ldots\) from \(S\) such that \(a_k\) minimizes the
\(p\)-adic valuation of \(\prod_{i=0}^{k-1} (a_k - a_i)\). - Define
\(v_k(S, p)\) to be the \(p\)-adic valuation of this minimal product. -
Then \(k!_S = \prod_p p^{v_k(S, p)}\).

\subsection{10.2 Properties}\label{ch2014manjulbhargava:properties}

The Bhargava factorial has the notable property that many classical
theorems involving the ordinary factorial extend to this generalized
setting:

\begin{itemize}
\tightlist
\item
 \textbf{Polynomial values:} For any primitive integer polynomial \(f\)
 of degree \(k\), the GCD of \(\{f(a) : a \in S\}\) divides \(k!_S\).
\item
 \textbf{Binomial coefficients:} The generalized binomial coefficients
 \(\binom{m+n}{n}_S = (m+n)!_S / (m!_S \cdot n!_S)\) are integers.
\item
 \textbf{Difference products:} The product of pairwise differences of
 any \(k+1\) elements of \(S\) is divisible by
 \(0!_S \cdot 1!_S \cdots k!_S\).
\end{itemize}

\subsection{10.3 Examples}\label{ch2014manjulbhargava:examples}

\begin{itemize}
\tightlist
\item
 If \(S = \mathbb{Z}\), then \(k!_S = k!\) (the ordinary factorial).
\item
 If \(S = 2\mathbb{Z}\) (even integers), then \(k!_S = 2^k k!\).
\item
 If \(S\) is the set of primes, the Bhargava factorial grows very
 rapidly, reflecting the sparseness of primes.
\item
 If \(S\) is the set of squares, then \(k!_S = (2k)!/2\).
\end{itemize}

\subsection{10.4 Applications}\label{ch2014manjulbhargava:applications}

The Bhargava factorial has found applications in: - \textbf{Fixed
divisors of polynomials:} Understanding when a polynomial takes values
divisible by a fixed integer. - \textbf{Integer-valued polynomials:}
Characterizing polynomials that map \(S\) to \(\mathbb{Z}\). -
\textbf{Ring theory:} Generalizing the factorial to arbitrary Dedekind
domains.

\begin{center}\rule{0.5\linewidth}{0.5pt}\end{center}

\section{11. The Average Rank of Elliptic
Curves}\label{ch2014manjulbhargava:the-average-rank-of-elliptic-curves}

\subsection{11.1 The Problem}\label{ch2014manjulbhargava:the-problem-1}

An elliptic curve over \(\mathbb{Q}\) is a smooth projective curve of
genus 1 with a rational point, typically given by a Weierstrass
equation:

\[
E: y^2 = x^3 + Ax + B
\]

The set of rational points \(E(\mathbb{Q})\) forms a finitely generated
abelian group (Mordell--Weil theorem):

\[
E(\mathbb{Q}) \cong \mathbb{Z}^r \oplus E(\mathbb{Q})_{\text{tors}}
\]

where \(r\) is the rank of \(E\) --- one of the important
invariants of an elliptic curve. Despite decades of study, the
distribution of ranks among elliptic curves was almost completely
unknown.

\subsection{11.2 The Minimalist
Conjecture}\label{ch2014manjulbhargava:the-minimalist-conjecture}

The Goldfeld--Katz--Sarnak ``minimalist conjecture'' predicts that: -
50\% of elliptic curves have rank 0 - 50\% have rank 1 - 0\% have rank
\(\geq 2\)

In particular, the average rank should be \(1/2\). But before Bhargava
and Shankar's work, it was not even known that the average rank is
finite.

\subsection{11.3 Bhargava and Shankar's
Strategy}\label{ch2014manjulbhargava:bhargava-and-shankars-strategy}

Bhargava, together with his student Arul Shankar, developed a powerful
new approach to the problem. The key idea is to study the Selmer groups
of elliptic curves. For a prime \(p\), the \(p\)-Selmer group
\(\mathrm{Sel}_p(E)\) sits in an exact sequence:

\[
0 \to E(\mathbb{Q})/pE(\mathbb{Q}) \to \mathrm{Sel}_p(E) \to \Sh(E)[p] \to 0
\]

where \(\Sh(E)\) is the Tate--Shafarevich group. The size of
\(\mathrm{Sel}_p(E)\) bounds the rank:

\[
\operatorname{rank}(E) \le \dim_{\mathbb{F}_p} \mathrm{Sel}_p(E)
\]

Thus, to bound the average rank, it suffices to bound the average size
of \(\mathrm{Sel}_p(E)\).

\subsection{11.4 The Central Results}\label{ch2014manjulbhargava:the-central-results}

In a series of landmark papers, Bhargava and Shankar proved:

\begin{itemize}
\tightlist
\item
 \textbf{Average size of \(\mathrm{Sel}_2(E)\):} 3 (2010)
\item
 \textbf{Average size of \(\mathrm{Sel}_3(E)\):} 4 (2010)
\item
 \textbf{Average size of \(\mathrm{Sel}_4(E)\):} 7 (2013)
\item
 \textbf{Average size of \(\mathrm{Sel}_5(E)\):} 6 (2013)
\end{itemize}

From these results, they deduced: - The average rank of all elliptic
curves over \(\mathbb{Q}\) is less than 0.885. - At least 83.75\% of
elliptic curves have rank 0 or 1. - At least 20.62\% of elliptic curves
have rank 0. - A positive proportion of elliptic curves have rank 1
(conditional on finiteness of Tate--Shafarevich groups).

\begin{center}\rule{0.5\linewidth}{0.5pt}\end{center}

\section{12. The Bhargava--Shankar Program on Selmer
Groups}\label{ch2014manjulbhargava:the-bhargavashankar-program-on-selmer-groups}

\subsection{12.1 The 2-Selmer Group}\label{ch2014manjulbhargava:the-2-selmer-group}

The first key result (2010) was the determination of the average size of
the 2-Selmer group. Bhargava and Shankar proved that when elliptic
curves over \(\mathbb{Q}\) are ordered by height, the average size of
\(\mathrm{Sel}_2(E)\) is exactly 3.

This implies that the limsup of the average rank is at most 1.5, since
\(\dim_{\mathbb{F}_2} \mathrm{Sel}_2(E) = \log_2 |\mathrm{Sel}_2(E)|\).

\subsection{12.2 The 3-Selmer Group}\label{ch2014manjulbhargava:the-3-selmer-group}

In a second paper (also 2010), they proved that the average size of
\(\mathrm{Sel}_3(E)\) is exactly 4. This gives a stronger bound: the
average rank is at most 1.17 (since \(\log_3 4 < 1.17\)).

More importantly, combining the 3-Selmer result with a parity theorem of
Dokchitser and Dokchitser, they proved that a positive proportion of
elliptic curves have rank 0.

\subsection{12.3 The 4-Selmer and 5-Selmer
Groups}\label{ch2014manjulbhargava:the-4-selmer-and-5-selmer-groups}

The 4-Selmer and 5-Selmer results (2013) were technically much more
demanding. The average sizes of 7 and 6 respectively implied the
sharpest bounds:

\begin{itemize}
\tightlist
\item
 Average rank \(<\) 0.885
\item
 At least 83.75\% have rank 0 or 1
\end{itemize}

These results represented the first definitive statements about the
distribution of ranks of elliptic curves.

\subsection{12.4 The Method: Counting
Orbits}\label{ch2014manjulbhargava:the-method-counting-orbits}

The common thread across all the Selmer group results is a powerful
technique for counting integral orbits in representations of algebraic
groups. Bhargava and Shankar showed that elements of the \(p\)-Selmer
group of an elliptic curve correspond to orbits of certain group actions
on spaces of forms of degree \(p\). By counting these orbits using
geometry-of-numbers methods, they obtained asymptotics for the number of
Selmer elements.

This method is a tour de force of analytic number theory: it combines
algebraic parametrizations, geometry-of-numbers estimates, and careful
analysis of local conditions to produce exact averages.

\begin{center}\rule{0.5\linewidth}{0.5pt}\end{center}

\section{13. The 2-Selmer Group and Binary Quartic
Forms}\label{ch2014manjulbhargava:the-2-selmer-group-and-binary-quartic-forms}

\subsection{13.1 Binary Quartic Forms}\label{ch2014manjulbhargava:binary-quartic-forms}

A binary quartic form is a homogeneous polynomial of degree 4 in two
variables:

\[
f(x, y) = a x^4 + b x^3 y + c x^2 y^2 + d x y^3 + e y^4
\]

Bhargava and Shankar showed that 2-Selmer elements of elliptic curves
correspond to \(\mathrm{GL}_2(\mathbb{Z})\)-orbits of binary quartic
forms satisfying certain local conditions.

\subsection{13.2 Counting Orbits}\label{ch2014manjulbhargava:counting-orbits}

The main technical achievement was to count the number of
\(\mathrm{GL}_2(\mathbb{Z})\)-orbits of binary quartic forms with
bounded invariants. This required extending the classical methods of
Gauss (for quadratic forms) and Davenport (for cubic forms) to the
quartic case.

The counting is complicated by the fact that the space of binary quartic
forms has a larger automorphism group and more complicated invariant
theory than the quadratic or cubic cases.

\subsection{13.3 The Average of 3}\label{ch2014manjulbhargava:the-average-of-3}

After counting orbits and analyzing the local conditions that
distinguish 2-Selmer elements, Bhargava and Shankar proved:

\[
\lim_{H \to \infty} \frac{\sum_{E: h(E) < H} |\mathrm{Sel}_2(E)|}{\#\{E: h(E) < H\}} = 3
\]

where \(h(E)\) is the height of the elliptic curve \(E\).

\begin{center}\rule{0.5\linewidth}{0.5pt}\end{center}

\section{14. Ternary Cubic Forms and the 3-Selmer
Group}\label{ch2014manjulbhargava:ternary-cubic-forms-and-the-3-selmer-group}

\subsection{14.1 Ternary Cubic Forms}\label{ch2014manjulbhargava:ternary-cubic-forms}

A ternary cubic form is a homogeneous polynomial of degree 3 in three
variables:

\[
F(x, y, z) = \sum_{i+j+k=3} a_{ijk} x^i y^j z^k
\]

The space of ternary cubic forms has 20 coefficients and carries an
action of \(\mathrm{SL}_3(\mathbb{Z})\).

\subsection{14.2 Connection to 3-Selmer
Groups}\label{ch2014manjulbhargava:connection-to-3-selmer-groups}

Bhargava and Shankar proved that 3-Selmer elements correspond to
\(\mathrm{SL}_3(\mathbb{Z})\)-orbits of ternary cubic forms. This is a
deeper and more subtle parametrization than the binary quartic case.

\subsection{14.3 Counting and the Average of
4}\label{ch2014manjulbhargava:counting-and-the-average-of-4}

Using geometry-of-numbers methods adapted to the 20-dimensional space of
ternary cubic forms, Bhargava and Shankar proved:

\[
\lim_{H \to \infty} \frac{\sum_{E: h(E) < H} |\mathrm{Sel}_3(E)|}{\#\{E: h(E) < H\}} = 4
\]

This result gave the first proof that a positive proportion of elliptic
curves have rank 0.

\begin{center}\rule{0.5\linewidth}{0.5pt}\end{center}

\section{15. The 4-Selmer and 5-Selmer
Groups}\label{ch2014manjulbhargava:the-4-selmer-and-5-selmer-groups-1}

\subsection{15.1 The 4-Selmer Group}\label{ch2014manjulbhargava:the-4-selmer-group}

For the 4-Selmer group, Bhargava and Shankar developed a more indirect
approach. Instead of parametrizing 4-Selmer elements directly (which
would involve 8-coverings of elliptic curves), they used a relation
between 2-Selmer and 4-Selmer groups.

Specifically, there is an exact sequence:

\[
0 \to \mathrm{Sel}_2(E) \to \mathrm{Sel}_4(E) \to \mathrm{Sel}_2(E)[2]
\]

Using this sequence together with their work on \(\mathrm{Sel}_2(E)\),
they proved that the average size of \(\mathrm{Sel}_4(E)\) is 7.

\subsection{15.2 The 5-Selmer Group}\label{ch2014manjulbhargava:the-5-selmer-group}

The 5-Selmer result was the most technically demanding. Bhargava and
Shankar developed a parametrization of 5-Selmer elements using spaces of
5×5 skew-symmetric matrices, building on earlier work of Bhargava on
quintic rings.

The proof involved: - A parametrization of 5-coverings of elliptic
curves - A geometry-of-numbers count in a 40-dimensional space - A
delicate analysis of local conditions and reducible orbits

The result --- that the average size of \(\mathrm{Sel}_5(E)\) is 6 ---
was a landmark in arithmetic statistics.

\subsection{15.3 The Full Sequence of
Averages}\label{ch2014manjulbhargava:the-full-sequence-of-averages}

The sequence of averages --- 3, 4, 7, 6 for \(p = 2, 3, 4, 5\) --- has a
striking structure. These numbers are not arbitrary: they arise from the
representation theory of the algebraic groups involved and from the
local mass formulas for Selmer groups. The fact that the averages are
exactly these integers (not just asymptotic bounds) reflects the deep
arithmetic structure underlying the theory.

\begin{center}\rule{0.5\linewidth}{0.5pt}\end{center}

\section{16. Consequences for the BSD Conjecture and Rank
Distribution}\label{ch2014manjulbhargava:consequences-for-the-bsd-conjecture-and-rank-distribution}

\subsection{16.1 The Birch and Swinnerton-Dyer
Conjecture}\label{ch2014manjulbhargava:the-birch-and-swinnerton-dyer-conjecture}

The Birch and Swinnerton-Dyer (BSD) conjecture, one of the seven Clay
Millennium Prize Problems, asserts that the rank of an elliptic curve
equals the order of vanishing of its \(L\)-function at \(s = 1\).
Bhargava and Shankar's results on Selmer groups have deep implications
for BSD.

\subsection{16.2 Analytic Rank 0}\label{ch2014manjulbhargava:analytic-rank-0}

Combining their 3-Selmer results with a theorem of Skinner and Urban on
the vanishing of \(L\)-functions, Bhargava and Skinner proved that a
positive proportion of elliptic curves have analytic rank 0 (i.e.,
non-vanishing \(L\)-function at \(s = 1\)). This implies, by the theorem
of Kolyvagin and the work of Skinner--Urban, that a positive proportion
of elliptic curves satisfy the BSD conjecture.

\subsection{16.3 Rank Distribution}\label{ch2014manjulbhargava:rank-distribution}

The combined results of Bhargava and Shankar give the following picture
of the rank distribution:

\begin{itemize}
\tightlist
\item
 At least 20.62\% of elliptic curves have rank 0
\item
 At least 83.75\% have rank 0 or 1
\item
 The average rank is less than 0.885
\item
 At most 16.25\% have rank \(\ge 2\)
\end{itemize}

These were the first unconditional results on the distribution of ranks
of elliptic curves. They show that the minimalist conjecture --- though
still unproven --- is consistent with all available evidence.

\subsection{16.4 Further Refinements}\label{ch2014manjulbhargava:further-refinements}

Subsequent work has refined these bounds: - Bhargava and Skinner (2014)
proved that a positive proportion of elliptic curves have rank exactly 1
(conditional on finiteness of Tate--Shafarevich groups). - Combining
with Dokchitser's parity results, one can show that among curves with
root number +1, a positive proportion have rank 0.

\begin{center}\rule{0.5\linewidth}{0.5pt}\end{center}

\section{17. Universal Quadratic Forms and the
290-Theorem}\label{ch2014manjulbhargava:universal-quadratic-forms-and-the-290-theorem}

\subsection{17.1 The Problem}\label{ch2014manjulbhargava:the-problem-2}

A quadratic form \(Q(x_1, \ldots, x_n)\) with integer coefficients is
called \textbf{universal} if it represents every positive integer.
Lagrange's four-square theorem (1770) states that
\(x^2 + y^2 + z^2 + w^2\) is universal.

A fundamental problem: which quadratic forms are universal? The
15-theorem of Conway and Schneeberger (1993) showed that if a positive
definite quadratic form represents the numbers 1 through 15, then it
represents all positive integers. However, this theorem applied only to
forms with integer matrix (i.e., ``classic'' forms).

\subsection{17.2 The 290-Theorem}\label{ch2014manjulbhargava:the-290-theorem}

In joint work with Jonathan Hanke, Bhargava proved the 290-Theorem: a
positive definite quadratic form with integer coefficients (not
necessarily with integer matrix) is universal if and only if it
represents every integer up to 290.

This was a dramatic extension of the Conway--Schneeberger theorem. The
proof involved: - An analysis of the representation properties of all
quadratic forms - A classification of the ``exceptional'' forms that
fail to represent certain integers - Extensive computer calculations to
verify the finite list of critical integers

\subsection{17.3 Impact}\label{ch2014manjulbhargava:impact-1}

The 290-Theorem has become a central result in the arithmetic theory of
quadratic forms. It provides a finite criterion for universality,
reducing a question about infinite sets to a finite check. The methods
developed in the proof have been applied to many related problems.

\begin{center}\rule{0.5\linewidth}{0.5pt}\end{center}

\section{18. Further Contributions}\label{ch2014manjulbhargava:further-contributions}

\subsection{18.1 Hyperelliptic Curves}\label{ch2014manjulbhargava:hyperelliptic-curves}

In joint work with Bjorn Poonen and Michael Stoll, Bhargava proved that
most hyperelliptic curves over \(\mathbb{Q}\) have no rational points.
More precisely, for odd degree hyperelliptic curves:

\[
C: y^2 = f(x)
\]

where \(f\) has degree \(2g+1\), a positive proportion of these curves
(when ordered by height) have no rational points other than the point at
infinity.

This result solved a long-standing problem about the distribution of
rational points on curves of genus \(\ge 2\).

\subsection{18.2 Plane Cubics and the Hasse
Principle}\label{ch2014manjulbhargava:plane-cubics-and-the-hasse-principle}

Bhargava proved that a positive proportion of plane cubic curves fail
the Hasse principle. The Hasse principle --- that a curve has rational
points if and only if it has points over every completion of
\(\mathbb{Q}\) --- holds for quadrics (by the Hasse--Minkowski theorem)
but fails in general. Bhargava's work showed that it fails for a
positive proportion of plane cubics, providing the first quantitative
evidence for the failure of local-global principles.

\subsection{18.3 The Hasse Principle for Hyperelliptic
Curves}\label{ch2014manjulbhargava:the-hasse-principle-for-hyperelliptic-curves}

Extending his work on plane cubics, Bhargava proved that a positive
proportion of hyperelliptic curves of genus 1 also fail the Hasse
principle, confirming that the failure of local-global principles is a
widespread phenomenon.

\subsection{18.4 Factorials and Fixed
Divisors}\label{ch2014manjulbhargava:factorials-and-fixed-divisors}

Beyond the Bhargava factorial itself, Bhargava has studied the fixed
divisors of polynomials --- the largest integer dividing \(f(a)\) for
all \(a\) in a given set \(S\). His work systematized the theory of
integer-valued polynomials and connected it to the factorial function in
novel ways.

\subsection{18.5 Coregular
Representations}\label{ch2014manjulbhargava:coregular-representations}

Bhargava's work on coregular representations --- representations of
algebraic groups whose ring of invariants is a polynomial ring ---
provided a systematic framework for studying the orbits that arise in
his parametrizations. This framework has been influential in
representation theory and invariant theory.

\begin{center}\rule{0.5\linewidth}{0.5pt}\end{center}

\section{19. Broader Impact and
Influence}\label{ch2014manjulbhargava:broader-impact-and-influence}

\subsection{19.1 The Creation of Arithmetic
Statistics}\label{ch2014manjulbhargava:the-creation-of-arithmetic-statistics}

Bhargava's work has created a new subfield of number theory: arithmetic
statistics. Before his work, the distribution of number fields, ranks of
elliptic curves, and Selmer groups were studied primarily through
heuristics and conjectures. Bhargava's methods made it possible to prove
precise theorems about these distributions.

The field now includes: - Counting number fields with specified
properties - Distribution of class groups and Selmer groups - Average
behavior of arithmetic invariants - Statistical properties of rational
points on curves

\subsection{19.2 Influence on Young
Researchers}\label{ch2014manjulbhargava:influence-on-young-researchers}

Bhargava has had a significant impact on the next generation of number
theorists. His students and collaborators include Arul Shankar, Jonathan
Hanke, Piper Harron, and many others who have extended his methods to
new problems. The techniques he developed --- counting orbits via the
geometry of numbers, parametrizing algebraic structures by forms, and
analyzing Selmer groups via representation theory --- have become
standard tools in arithmetic geometry.

\subsection{19.3 Connections to Physics}\label{ch2014manjulbhargava:connections-to-physics}

Bhargava's higher composition laws have found unexpected applications in
theoretical physics. The ``charge cube'' of STU black holes in string
theory has exactly the same structure as Bhargava's cube, and the higher
composition laws describe the U-duality orbits of black hole solutions.
This connection, discovered by physicists working on supergravity and
M-theory, reveals a deep link between number theory and quantum gravity.

\subsection{19.4 Public Engagement}\label{ch2014manjulbhargava:public-engagement}

Bhargava is one of the most visible public faces of mathematics. He has
given numerous public lectures, appeared in documentaries, and served as
the first Distinguished Chair for Public Dissemination of Mathematics at
the National Museum of Mathematics. His engaging style --- blending
mathematics with magic tricks, games, and music --- has inspired
countless young people to pursue mathematics.

\begin{center}\rule{0.5\linewidth}{0.5pt}\end{center}

\section{20. Awards and Recognition}\label{ch2014manjulbhargava:awards-and-recognition}

\textbf{2019:} - Fellow of the Royal Society (FRS)

\textbf{2018:} - Distinguished Chair for Public Dissemination of
Mathematics, National Museum of Mathematics

\textbf{2015:} - Padma Bhushan, India's third-highest civilian award

\textbf{2014:} - Fields Medal (International Mathematical Union)

\textbf{2013:} - Elected to the U.S. National Academy of Sciences

\textbf{2012:} - Simons Investigator Award (inaugural class) - Infosys
Prize in Mathematics - Fellow, American Mathematical Society

\textbf{2011:} - Fermat Prize

\textbf{2008:} - Cole Prize in Number Theory (American Mathematical
Society)

\textbf{2005:} - Clay Research Award - SASTRA Ramanujan Prize -
Blumenthal Award for the Advancement of Research in Pure Mathematics

\textbf{2004:} - Hasse Prize (Mathematical Association of America)

\textbf{2003:} - Clay Research Fellowship

\textbf{2002:} - Popular Science ``Brilliant 10'' Award

\textbf{2000:} - Morgan Prize for Outstanding Undergraduate Research in
Mathematics

\begin{center}\rule{0.5\linewidth}{0.5pt}\end{center}

\section{21. Selected Major
Publications}\label{ch2014manjulbhargava:selected-major-publications}

\begin{enumerate}
\def\labelenumi{\arabic{enumi}.}
\item
 \textbf{``Higher Composition Laws I: A New View on Gauss Composition,
 and Quadratic Generalizations.''} \emph{Annals of Mathematics} 159(1),
 217--250, 2004. --- \emph{The foundational paper introducing the
 Bhargava cube and the first new composition laws.}
\item
 \textbf{``Higher Composition Laws II: On Cubic Analogues of Gauss
 Composition.''} \emph{Annals of Mathematics} 159(2), 865--886, 2004.
 --- \emph{Cubic composition laws and the 3×3×3 cube.}
\item
 \textbf{``Higher Composition Laws III: The Parametrization of Quartic
 Rings.''} \emph{Annals of Mathematics} 159(3), 1329--1360, 2004. ---
 \emph{Parametrization of quartic rings by pairs of ternary quadratic
 forms.}
\item
 \textbf{``Higher Composition Laws IV: The Parametrization of Quintic
 Rings.''} \emph{Annals of Mathematics} 167(1), 53--94, 2008. ---
 \emph{Parametrization of quintic rings by quadruples of alternating
 5×5 forms.}
\item
 \textbf{``The Density of Discriminants of Quartic Rings and Fields.''}
 \emph{Annals of Mathematics} 162(2), 1031--1063, 2005. ---
 \emph{Asymptotic count of quartic fields with bounded discriminant.}
\item
 \textbf{``The Density of Discriminants of Quintic Rings and Fields.''}
 \emph{Annals of Mathematics} 172(3), 1559--1591, 2010. ---
 \emph{Asymptotic count of quintic fields.}
\item
 \textbf{``The Factorial Function and Generalizations.''}
 \emph{American Mathematical Monthly} 107(9), 783--799, 2000. ---
 \emph{Introduction of the Bhargava factorial.}
\item
 \textbf{``Binary Quartic Forms Having Bounded Invariants, and the
 Boundedness of the Average Rank of Elliptic Curves''} (with A.
 Shankar). \emph{Annals of Mathematics} 181(1), 191--242, 2015. ---
 \emph{Average size of 2-Selmer groups is 3; boundedness of average
 rank.}
\item
 \textbf{``Ternary Cubic Forms Having Bounded Invariants, and the
 Existence of a Positive Proportion of Elliptic Curves Having Rank 0''}
 (with A. Shankar). \emph{Annals of Mathematics} 181(2), 587--621,
 2015. --- \emph{Average size of 3-Selmer groups is 4; positive
 proportion of rank 0.}
\item
 \textbf{``The Average Size of the 5-Selmer Group of Elliptic Curves is
 6, and the Average Rank is Less Than 1''} (with A. Shankar).
 arXiv:1312.7859, 2013. --- \emph{Average rank \textless{} 0.885; at
 least 83.75\% have rank 0 or 1.}
\item
 \textbf{``The Average Number of Elements in the 4-Selmer Groups of
 Elliptic Curves is 7''} (with A. Shankar). arXiv:1312.7333, 2013. ---
 \emph{Average size of 4-Selmer groups is 7.}
\item
 \textbf{``Universal Quadratic Forms and the 290-Theorem''} (with J.
 Hanke). In preparation. --- \emph{The 290-Theorem for universal
 quadratic forms.}
\item
 \textbf{``A Positive Proportion of Plane Cubics Fail the Hasse
 Principle.''} arXiv:1402.1131, 2014. --- \emph{Quantitative failure of
 local-global principle.}
\item
 \textbf{``Most Odd Degree Hyperelliptic Curves Have Only One Rational
 Point''} (with B. Poonen and M. Stoll). \emph{Annals of Mathematics}
 192(3), 2020. --- \emph{Distribution of rational points on
 hyperelliptic curves.}
\end{enumerate}

\begin{center}\rule{0.5\linewidth}{0.5pt}\end{center}

\section{22. Ongoing Work and Future
Directions}\label{ch2014manjulbhargava:ongoing-work-and-future-directions}

\subsection{22.1 Higher Degree Fields}\label{ch2014manjulbhargava:higher-degree-fields}

Bhargava's parametrizations and counting methods currently work up to
degree 5. Extending these results to degree 6 and beyond is a major open
problem. The obstacle is that the relevant parametrizations become
increasingly complex, and the geometry-of-numbers estimates become more
difficult in higher dimensions.

\subsection{\texorpdfstring{22.2 Average Ranks for
\(p \ge 7\)}{22.2 Average Ranks for p \textbackslash ge 7}}\label{ch2014manjulbhargava:average-ranks-for-p-ge-7}

The Bhargava--Shankar results cover \(p = 2, 3, 4, 5\). For \(p \ge 7\),
the parametrizations of Selmer groups become more complex, and the
counting problems are higher-dimensional. Determining the average size
of \(\mathrm{Sel}_p(E)\) for all \(p\) would give even sharper bounds on
the average rank and potentially prove the minimalist conjecture.

\subsection{22.3 The Mordell--Weil and Tate--Shafarevich
Groups}\label{ch2014manjulbhargava:the-mordellweil-and-tateshafarevich-groups}

The results on Selmer groups give information about the quotient
\(E(\mathbb{Q})/pE(\mathbb{Q})\) and the \(p\)-torsion of the
Tate--Shafarevich group. Understanding the distribution of ranks
themselves --- not just Selmer ranks --- is the ultimate goal. The
minimalist conjecture remains the central open problem.

\subsection{22.4 Local-to-Global
Principles}\label{ch2014manjulbhargava:local-to-global-principles}

Bhargava's work on the Hasse principle for plane cubics and
hyperelliptic curves has opened up the statistical study of
local-to-global principles. Future work includes understanding the
proportion of varieties that satisfy the Hasse principle, the
Brauer--Manin obstruction, and other refinements.

\subsection{22.5 Connections to Algebraic
Geometry}\label{ch2014manjulbhargava:connections-to-algebraic-geometry}

The higher composition laws have deep connections to the theory of
algebraic stacks and moduli spaces. Bhargava's orbit-counting methods
can be reinterpreted as counting rational points on certain moduli
stacks, connecting arithmetic statistics to algebraic geometry.

\begin{center}\rule{0.5\linewidth}{0.5pt}\end{center}

\section{23. Bhargava's Mathematical
Style}\label{ch2014manjulbhargava:bhargavas-mathematical-style}

\subsection{23.1 Problem Selection}\label{ch2014manjulbhargava:problem-selection}

Bhargava has a notable instinct for choosing problems that are both
central and tractable. His targets are typically classical problems that
have resisted progress for decades or centuries: the generalization of
Gauss composition, the density of discriminants of number fields, the
average rank of elliptic curves. Each of these problems was well-known
to be important, but no one knew how to approach them.

\subsection{23.2 Conceptual Clarity}\label{ch2014manjulbhargava:conceptual-clarity}

One of Bhargava's greatest strengths is his ability to find simple,
conceptual formulations of complex problems. The Bhargava cube is the
perfect example: a 200-year-old theory reduced to a Rubik's cube
pattern. This clarity extends to his papers, which are models of
exposition.

\subsection{23.3 Integration of Classical and
Modern}\label{ch2014manjulbhargava:integration-of-classical-and-modern}

Bhargava's work is notable for its integration of classical number
theory with modern algebraic techniques. He reads Gauss and Dirichlet as
carefully as he reads contemporary papers, and his work often shows that
classical ideas contain hidden depths that modern methods can unlock.

\subsection{23.4 Collaboration}\label{ch2014manjulbhargava:collaboration}

Bhargava is a highly collaborative mathematician. His important
work --- on elliptic curves --- was done with his student Arul Shankar,
and he has co-authored with a wide range of mathematicians including
Andrew Wiles, Benedict Gross, Christopher Skinner, Bjorn Poonen, Michael
Stoll, and many others. He is known for his generosity with ideas and
his mentorship of young mathematicians.

\begin{center}\rule{0.5\linewidth}{0.5pt}\end{center}

\section{24. Conclusion}\label{ch2014manjulbhargava:conclusion}

Manjul Bhargava is one of the most original and influential number
theorists of his generation. In a career that began with his PhD in
2001, he has: - Discovered 14 higher composition laws generalizing
Gauss's classical theory - Parametrized quartic and quintic rings,
solving problems that had been open for over a century - Determined the
asymptotic density of quartic and quintic fields - Invented the Bhargava
factorial, a fundamental new concept in number theory - Proved the
290-Theorem on universal quadratic forms - Determined the average size
of 2-, 3-, 4-, and 5-Selmer groups of elliptic curves - Proved that the
average rank of elliptic curves is bounded and less than 1 - Shown that
a positive proportion of elliptic curves have rank 0 and rank 1 - Proved
that most hyperelliptic curves have no rational points - Shown that a
positive proportion of plane cubics fail the Hasse principle

Each of these achievements would be sufficient for a distinguished
career; together, they represent one of the most notable bodies of
work in modern number theory.

The Fields Medal citation --- \emph{``for developing powerful new
methods in the geometry of numbers, which he applied to count rings of
small rank and to bound the average rank of elliptic curves''} ---
captures the essence of his contributions. His work has transformed our
understanding of the distribution of fundamental arithmetic objects,
from number fields to elliptic curves, and has created the new field of
arithmetic statistics. The methods he developed --- particularly the use
of integral orbit counting in the geometry of numbers --- will continue
to influence mathematics for generations to come.

\begin{center}\rule{0.5\linewidth}{0.5pt}\end{center}

\begin{center}\rule{0.5\linewidth}{0.5pt}\end{center}

\section*{Exercises}
\addcontentsline{toc}{section}{Exercises}

\begin{exercise} Compute the 3-class group of the field $\\mathbb{Q}(\\sqrt{-23})$ and verify Gauss's class number problem for this case. \end{exercise}

\begin{exercise} Show that the number of $\\mathbb{Z}$-isomorphism classes of quadratic rings of discriminant $D$ is finite. \end{exercise}

\begin{exercise} Verify Bhargava's formula: the number of $\\mathbb{Z}$-lattices in $\\mathbb{Q}^n$ of index $N$ is $\\prod_{p|N} \\frac{p^n-1}{p-1}$. \end{exercise}

\begin{exercise} Prove that the average number of 2-Selmer elements of elliptic curves over $\\mathbb{Q}$ is bounded (Bhargava-Shankar theorem). \end{exercise}

\begin{center}\rule{0.5\linewidth}{0.5pt}\end{center}

\section*{Summary}
\addcontentsline{toc}{section}{Summary}

Bhargava invented higher composition laws, providing a new framework for counting number fields and solving classical problems in number theory. His proof that the average rank of elliptic curves over $\mathbb{Q}$ is bounded (with Shankar) is a major advance toward the Birch and Swinnerton-Dyer conjecture, and his work on the geometric composition law connects number theory to algebraic geometry.

\section{25. References and Further
Reading}\label{ch2014manjulbhargava:references-and-further-reading}

\subsection{Official Sources}\label{ch2014manjulbhargava:official-sources}

\begin{itemize}
\tightlist
\item
 IMU Fields Medal 2014 Announcement:
 https://www.mathunion.org/imu-awards/fields-medal/fields-medals-2014
\item
 Fields Medal Citation for Manjul Bhargava:
 https://www.mathunion.org/fileadmin/IMU/Prizes/Fields/2014/news\_release\_bhargava.pdf
\item
 Soundararajan, K. ``The Work of Manjul Bhargava.'' ICM 2014 Laudatio.
\end{itemize}

\subsection{Biographical and Overview
Articles}\label{ch2014manjulbhargava:biographical-and-overview-articles}

\begin{itemize}
\tightlist
\item
 Klarreich, E. ``The Musical, Magical Number Theorist.'' \emph{Quanta
 Magazine}, August 12, 2014.
\item
 Krieger, L. M. ``Manjul Bhargava: The Number Theorist Who Would Be
 King.'' \emph{Nature}, August 2014.
\item
 Princeton University News. ``Bhargava Receives Fields Medal.'' August
 12, 2014.
\item
 Lamb, E. ``The Unlikely Mathematician Who Modernized Number Theory.''
 \emph{Scientific American}, 2014.
\end{itemize}

\subsection{Key Technical References}\label{ch2014manjulbhargava:key-technical-references}

\begin{itemize}
\tightlist
\item
 Bhargava, M. ``Higher Composition Laws I--IV.'' \emph{Annals of
 Mathematics}, 2004--2008.
\item
 Bhargava, M. ``The Density of Discriminants of Quartic Rings and
 Fields.'' \emph{Annals of Mathematics} 162 (2005), 1031--1063.
\item
 Bhargava, M. ``The Density of Discriminants of Quintic Rings and
 Fields.'' \emph{Annals of Mathematics} 172 (2010), 1559--1591.
\item
 Bhargava, M. and Shankar, A. ``Binary Quartic Forms Having Bounded
 Invariants, and the Boundedness of the Average Rank of Elliptic
 Curves.'' \emph{Annals of Mathematics} 181 (2015), 191--242.
\item
 Bhargava, M. and Shankar, A. ``Ternary Cubic Forms Having Bounded
 Invariants, and the Existence of a Positive Proportion of Elliptic
 Curves Having Rank 0.'' \emph{Annals of Mathematics} 181 (2015),
 587--621.
\item
 Bhargava, M. and Shankar, A. ``The Average Size of the 5-Selmer Group
 of Elliptic Curves is 6.'' arXiv:1312.7859.
\item
 Bhargava, M. and Hanke, J. ``Universal Quadratic Forms and the
 290-Theorem.'' In preparation.
\item
 Bhargava, M., Poonen, B., and Stoll, M. ``Most Odd Degree
 Hyperelliptic Curves Have Only One Rational Point.'' \emph{Annals of
 Mathematics} 192 (2020).
\end{itemize}

\begin{center}\rule{0.5\linewidth}{0.5pt}\end{center}

\section*{See Also}
\addcontentsline{toc}{section}{See Also}

This chapter relates to the following work by other Fields Medalists:
\begin{itemize}
  \item \textbf{\hyperref[ch1950atleselberg:atle-selberg-fields-medal-1950-a-comprehensive-report]{See Chapter on related work}} --- Selberg's trace formula relates to Bhargava.
  \item \textbf{\hyperref[ch1974enricobombieri:enrico-bombieri-fields-medal-1974-a-comprehensive-report-on-his-work-and-impact]{See Chapter on related work}} --- Bombieri's sieve methods relate to Bhargava.
  \item \textbf{\hyperref[ch2006terencetao:terence-tao-fields-medal-2006-a-comprehensive-report-on-his-work-and-impact]{See Chapter on related work}} --- Tao's analytic number theory relates to Bhargava.
\end{itemize}

